\summarypart{FOUR}{The Christian Middle Ages}{590--1517}
\begin{partoverview}
Europe owes more to the Christian faith than most people realize. When the barbarians destroyed the
Roman Empire in the West, it was the church that assembled a new order called Europe, taking the lead
in the rule of law, the pursuit of knowledge and the expressions of culture. The idea beneath it all
was \emph{Christendom}, the union of empire and church. It began under Charlemagne in the eighth
century, but the popes slowly gathered more and more power, until Innocent~III (1198--1216) taught
Europe to think of the popes as rulers of the world. The later centuries saw the papacy corrupted by
that power, and reformers, increasingly militant, crying out for change.
\end{partoverview}
\partchapters{%
  \pch{17}{God's Consul}{Gregory the Great}
  \pch{18}{The Search for Unity}{Charlemagne and Christendom}
  \pch{19}{Lifted in a Mystic Manner}{The Papacy and the Crusader}
  \pch{20}{The Nectar of Learning}{Scholasticism}
  \pch{21}{A Song to Lady Poverty}{The Apostolic Lifestyle}
  \pch{22}{Sleeping Men and the Law of Necessity}{The Decline of the Papacy}
  \pch{23}{Judgment in the Process of Time}{Wyclif and Hus}}
